\chapter{The page}
\label{chap:page}

Medorama is used through one web page, served by \cmd{medorama serve} (by default at
\url{http://127.0.0.1:8090}). This chapter is for every user. It describes what is common to the
whole page: signing in, the header and its tabs, the addresses, the light and dark themes, the
Privacy Mode switch, the notices, the keyboard and the History tab. The following chapters take
the tabs one by one.

\section{Signing in}
\label{sec:page-signing-in}

How the page signs users in depends on the catalog's sign-in method, which an administrator
chooses with the setting \code{sign\_in} (\cref{chap:sign-in}). The page asks the server which
method is in use when it shows the sign-in screen.

\paragraph{With Medorama's passwords} (the default), the page opens on a sign-in form with the
fields \ui{User name} (the account's name, for example \code{admin}) and \ui{Password}, and the
button \ui{Sign in}. An administrator creates the accounts, in the \ui{Admin} tab or with
\cmd{medorama user add} (\cref{chap:users-privacy}); \cmd{medorama init} creates the first one,
\code{admin}.
\begin{itemize}
\item A wrong user name or password is reported above the form. After 5 failed sign-ins, a user
  name is locked for a minute (``too many failed attempts'').
\item A sign-in lasts \code{session\_hours} (12 hours by default), and ends sooner when it is not
  used for \code{session\_idle\_minutes} (60 minutes by default). The page's own checks for new
  data do not count as use. When the session has ended, the next request brings the sign-in form
  back.
\item Changing a password or disabling a user ends that user's sessions.
\end{itemize}

\paragraph{With the institution's sign-in} (\code{sign\_in = "proxy"}), an authenticating proxy
in front of Medorama signs the user in before the page loads, for example through single sign-on
and multi-factor authentication. The page shows no form. When the proxy has not signed the user
in, or the proxy's session has ended, the page says so and offers \ui{Sign in}, which loads the
page again through the proxy. When the proxy names someone for whom Medorama has no account, the
page shows the server's reason: ``Your institution's sign-in names you \textit{name}, and
Medorama has no account of that name: ask an administrator.''

\paragraph{The user menu.} The round button at the right of the header shows the user's
initials. It opens a menu with the user's name, role and, for users with the PHI permission, the
mark \ui{PHI}, and these commands:
\begin{description}
\item[\ui{Change password}] asks for the current password, then twice for the new one (at least
  10 characters); the user then signs in again.
\item[\ui{Sign out}] ends this session. With the institution's sign-in, it leads to the proxy's
  own sign-out address (the setting \code{proxy\_logout\_url}), or to the signed-out screen when
  none is set.
\item[\ui{Sign out everywhere}] ends the user's sessions on every computer and browser.
\end{description}
With the institution's sign-in, Medorama's passwords are not used, so the menu hides
\ui{Change password} and \ui{Sign out everywhere}, and \ui{Admin} $\rightarrow$ \ui{Users}
shows no password field and no \ui{Reset password}.

\section{The header}
\label{sec:page-header}

The header, at the top of every screen, holds from left to right the name Medorama, the
\ui{$\langle$} and \ui{$\rangle$} buttons (\ui{Back} and \ui{Forward}: the browser's history of
this page), the tabs, the \ui{Privacy Mode} switch, the theme button and the user menu.

The tabs are in three groups (\cref{tab:page-tabs}):
\begin{itemize}
\item \ui{Cohorts}, \ui{Patient Catalog} and \ui{Databases} are three ways to the patients. The
  page opens on \ui{Cohorts}.
\item After a separator, \ui{ECG}, \ui{PACS} and \ui{VitalStream} are the viewers.
\item At the right, in a quieter style, \ui{History}, \ui{Actions}, \ui{Jobs}, \ui{Review},
  \ui{Import} and \ui{Admin} hold the work. Some of them are shown only to the roles that use
  them.
\end{itemize}

\begin{xltabular}{\linewidth}{@{}P{2.6cm}P{2.4cm}L@{}}
\caption{The tabs of the header.}
\label{tab:page-tabs}\\
\toprule
{\tabhead Tab} & {\tabhead Shown to} & {\tabhead Shows} \\
\midrule
\endfirsthead
\toprule
{\tabhead Tab} & {\tabhead Shown to} & {\tabhead Shows} \\
\midrule
\endhead
\bottomrule
\endfoot
\ui{Cohorts} & Everyone & Find patients across every database by criteria, with why each
  matched; save them as named cohorts; edit, clone, merge, refine, refresh, share and export
  them (\cref{chap:cohorts}). \\
\ui{Patient Catalog} & Everyone & Every patient, latest encounter first, with the kinds of data
  each has; a patient's page; the table of encounters (\cref{chap:patients}). \\
\ui{Databases} & Everyone & Every database in use, each unfolding into its rows as imported
  (\cref{chap:databases-tab}). \\
\ui{ECG}, \ui{PACS} & Everyone & The ECG and PACS viewers, at the file a patient's page or a
  database's row sent them. Dimmed until something is opened in them
  (\cref{chap:viewers-in-medorama}). \\
\ui{VitalStream} & Everyone & The waveform viewer the catalog runs. Grayed out, with the reason
  in its tooltip, when no viewer is linked to the catalog; absent when the catalog links to a
  viewer running elsewhere, which opens in a browser tab of its own. \\
\ui{History} & Everyone & The encounters the user opened in the encounter panel
  (\cref{sec:page-history}). \\
\ui{Actions} & Analysts, admins & Run actions on recordings (\cref{chap:actions}). \\
\ui{Jobs} & Everyone & The queue of jobs; a number by the tab counts the running and queued
  ones (\cref{chap:actions}). \\
\ui{Review} & Admins & Rows held back, possible duplicate patients, recordings whose files
  disagree with the databases; a number by the tab counts what is open
  (\cref{chap:importing}). \\
\ui{Import} & Admins & Add databases, the connected databases, import formats, sources and
  imports, the recordings; while an import runs, a bar by the tab shows its progress
  (\cref{chap:importing}, \cref{chap:connecting}, \cref{chap:recordings}). \\
\ui{Admin} & Admins & Users, settings, add-ons, dates and times, Privacy Mode's settings, the
  audit log (\cref{chap:administration}, \cref{chap:users-privacy}). \\
\end{xltabular}

\paragraph{Roles.} A \code{viewer} searches and looks, keeps a history and makes cohorts. An
\code{analyst} may also run and cancel actions. An \code{admin} may also add and remove data,
review identities, manage users and settings, and read the audit log. The PHI permission, given
per user, lets a user leave Privacy Mode. Roles and the permission are described in
\cref{chap:users-privacy}.

\section{Addresses}
\label{sec:page-addresses}

Every screen has its own address, after the \code{\#} of the page's URL
(\cref{tab:page-addresses}). The browser's back and forward buttons, and the header's
\ui{$\langle$} and \ui{$\rangle$}, move between the patient page and the viewers. A screen can be
reloaded, bookmarked, or sent to someone who may see it. An address that the page does not know
opens \ui{Cohorts}.

\begin{xltabular}{\linewidth}{@{}P{6.2cm}L@{}}
\caption{The addresses of the page's screens.}
\label{tab:page-addresses}\\
\toprule
{\tabhead Address} & {\tabhead Screen} \\
\midrule
\endfirsthead
\toprule
{\tabhead Address} & {\tabhead Screen} \\
\midrule
\endhead
\bottomrule
\endfoot
\code{\#/cohorts} & A new cohort: the criteria builder \\
\code{\#/cohorts/7}, \code{\#/cohorts/7/edit} & A cohort; its editing \\
\code{\#/cohorts/new?within=7} & A new subcohort of cohort 7 \\
\code{\#/patients?q=\dots\&has=ecg,!icu\&cohort=7} & The Patient Catalog, with its search,
  its filters and a cohort's members \\
\code{\#/encounters} & The table of encounters \\
\code{\#/patient/12?item=\dots\&enc=\dots\&q=\dots\&mode=\dots} & A patient's page, with the
  entry chosen, the encounter shown, the report's search and its mode \\
\code{\#/databases}, \code{\#/databases/12} & The databases; one on a page of its own \\
\code{\#/viewer?path=\dots\&t0=\dots\&t1=\dots} & VitalStream, at a recording and a time span \\
\code{\#/ecg?src=\dots}, \code{\#/pacs?src=\dots}, \code{\#/pacs?list=\dots} & The ECG and PACS
  viewers, at a file or a folder \\
\code{\#/history}, \code{\#/actions}, \code{\#/jobs} & History, Actions (\code{?cohort=7}: on a
  cohort's members), Jobs \\
\code{\#/review}, \code{\#/import}, \code{\#/recordings}, \code{\#/admin} & The admins' tabs \\
\end{xltabular}

\paragraph{Changes show everywhere at once.} Every open page asks the server for the catalog's
data version every 5 seconds, and at once after it has changed something itself. An import or a
removal, made here or in anyone else's browser, changes that version. The patient pages, the
lists and the tables then redraw, and a viewer showing a file that is gone closes with a note
(\cref{sec:page-notices}).

\section{Light and dark}
\label{sec:page-themes}

The page and every viewer have a light and a dark theme. The page starts in the system's setting,
that is, the light or dark appearance of the operating system or the browser. The theme button at
the right of the header (a half-filled circle, titled \ui{Light or dark}) switches between the
two themes. The choice is remembered in this browser, for this page, and then applies instead of
the system's setting.

The three viewers, in their tabs, follow the page's theme: the page sets it in their address
(\code{theme=light} or \code{theme=dark}) when a tab opens and when the theme changes. A viewer
applies it without remembering it as its own choice; VitalStream's own theme button, in its
toolbar, still switches its theme (\cref{chap:vitalstream}).

The screenshots of this manual show the dark theme, except those of the ECG viewer, which is
shown in both themes (\cref{chap:ecg-viewer}).

\screenshot{patient-icu.png}{The dark theme: a patient's page with an ICU recording at a
glance in the detail. Synthetic data.}{fig:page-dark}

Both themes use the same colors, each tuned to keep its contrast. The hues come from palettes
designed for color-vision deficiencies (Okabe-Ito, Paul Tol). Red is reserved for alarms. Amber
and orange mark abnormal and critical values and errors, and, in Privacy Mode, the warning
triangle on files that may show identifiers; blue marks a low value. No category of data uses
these colors.

\section{Dates and times}
\label{sec:page-dates}

Every page writes dates and times in the catalog's time zone (the setting \code{timezone}), in
the order and on the clock that an administrator chooses for everyone in \ui{Admin}
(\ui{Dates and times}): 2024-01-31, 01/31/2024, 31/01/2024 or 31.01.2024, with a 24-hour or
12-hour clock. Exports and
the API keep ISO 8601. In Privacy Mode, the dates shown are moved back
(\cref{sec:page-privacy}).

\section{The Privacy Mode switch}
\label{sec:page-privacy}

The \ui{Privacy Mode} button of the header shows and changes the session's mode:
\begin{itemize}
\item \textbf{On} (the button filled): no identifier is shown, and dates are moved back. Its
  tooltip reads ``Privacy Mode is on: no identifier is shown, dates are moved back. Click to show
  identifiers.''
\item \textbf{Off} (the button plain): identifiers are shown. The tooltip reads ``Privacy Mode
  is off: identifiers are shown. Click to hide them.''
\item \textbf{Locked on}, for users without the PHI permission: a click only says ``Privacy Mode
  is always on for your account''.
\end{itemize}
Administrators with the PHI permission start with Privacy Mode off; everyone else starts with it
on. A click asks the server to switch the session's mode, and the server writes the change to
the audit log. The page then reloads, so that nothing identified stays on the screen. In a narrow
window, the button shows only its symbol.

\begin{phibox}
The switch only asks; the server decides. In Privacy Mode the server sends no identifier at
all: placeholders (\code{[Patient Name]}, \code{[MRN]}), pseudonyms (\code{MO-XXXX-XXXX}),
moved dates and de-identified files take their place before anything leaves it. The page cannot
show what it is not sent. When a viewer's tab shows a picture, a PDF, a video or a DICOM file,
whose pixels or pages are sent as they are, the tab is marked with an amber warning triangle,
because such a file may show burned-in identifiers. What Privacy Mode withholds, and how dates
move, is explained in \cref{chap:protecting}.
\end{phibox}

\section{Notices}
\label{sec:page-notices}

\begin{itemize}
\item \textbf{Messages} appear at the bottom of the page for a few seconds. They say what an
  action did (``Saved'', ``Indexing queued'', ``3 patients added to \dots''); errors appear in
  amber, a little longer. A screen never fails silently: whatever goes wrong is reported there.
\item \textbf{Confirmations} come before every removal, saying first what would go
  (\cref{chap:databases-tab}, \cref{chap:cohorts}).
\item \textbf{A viewer whose file is gone} (its database was deleted) closes, and its tab says
  so: ``What this ECG tab showed was removed from the catalog (its database was deleted).''
\item \textbf{A patient no longer in the catalog} (its databases deleted, or merged into
  another): the patient's page says so, with a link to the Patient Catalog.
\item \textbf{A connected database that did not answer}: the patient's page marks it
  (``omop: not available'') and shows the other sources (\cref{chap:patients}).
\item \textbf{New data since a cohort's members were found}: the cohort's page says so, and
  offers to see who would join (\cref{chap:cohorts}).
\item \textbf{Counts by the tabs}: the jobs running and queued (\ui{Jobs}), what waits for
  review (\ui{Review}), the progress of the imports running (\ui{Import}).
\end{itemize}

\section{The keyboard}
\label{sec:page-keyboard}

\begin{table}[htbp]
\caption{Keys of the page.}
\label{tab:page-keys}
\begin{tabularx}{\linewidth}{@{}P{3.4cm}L@{}}
\toprule
{\tabhead Key} & {\tabhead Does} \\
\midrule
\key{/} & In the Patient Catalog and on a patient's page (when no field has the focus): puts the
  focus in the patient search. \\
\key{Enter} & In the Patient Catalog's search, when it finds one patient: opens their page. In
  the suggestions of a patient's page: opens the patient chosen. On an entry of the report, an
  encounter, a trend or a note: chooses it, as a click does. \\
\key{Space} & On an entry of the report, an encounter, a trend or a note: chooses it. \\
\key{$\uparrow$} \key{$\downarrow$} & In the suggestions of the search for another patient:
  move the choice. \\
\key{Esc} & Closes the suggestions, the encounter panel, and the \ui{New / Edit cohort}
  menu. \\
\key{Tab} & Moves through the page's controls; after a screen is drawn, the focus is on it, so
  that \key{Tab} goes on from there. \\
\bottomrule
\end{tabularx}
\end{table}

The browser's own keys for back and forward (for example \key{Alt}+\key{$\leftarrow$}) move in
the page's history, as the header's \ui{$\langle$} \ui{$\rangle$} do. The viewers have keys of
their own (\cref{part:viewers}).

\section{The History tab}
\label{sec:page-history}

\ui{History} lists the encounters that the user opened in the encounter panel, from the table
of encounters (\cref{chap:patients}), from \ui{Review}, or from the History itself; each opening
is a \code{GET encounters/\{id\}}. The list shows the most recent first, with the status of
their recordings, the encounter, the patient (and, for users with the PHI permission, the
name\phimark), the admission, when it was last opened, and how many times. Patient pages are not
recorded. Only the user sees their own history.

\begin{itemize}
\item A click on a row opens the encounter's panel again.
\item The $\times$ at the end of a row (\ui{Remove from your history}) removes that encounter.
\item \ui{Clear history} empties it, after a confirmation.
\end{itemize}
